\label{sec:axioms}
%==================================================================

Пока не доказано, что $\Shv{F}^{\#}$ --- пучок, все рассуждения о нём
--- лишь предположение. Проверим аксиомы.

\begin{rem}[алгебраическая структура]
На $\Shv{F}^{\#}(U)$ вводится поэлементная абелева структура (если
$\Shv{F}$ --- предпучок групп):
\[
 (s_1+s_2)(x)=s_1(x)+s_2(x),\qquad
 0(x)=0_{\Shv{F}_x}.
\]
Нулевой элемент --- непрерывное сечение, постоянно равное нулевому
ростку; оно непрерывно, потому что прообраз любого открытого ---
либо $\varnothing$, либо всё $U$. Остаётся проверить (F1) и (F2).
\end{rem}

\subsection{Аксиома (F1): единственность}
\textbf{Задача.} $U=\bigcup_i U_i$, $s',s''\in\Shv{F}^{\#}(U)$ и
$s'|_{U_i}=s''|_{U_i}$ для всех $i$. Доказать $s'=s''$.

\textbf{Доказательство.} Обозначим
\[
 E:=\{\,x\in U:\ s'(x)=s''(x)\,\}.
\]
Шаг 1: $E$ открыто. Действительно, пусть $s'(x)=s''(x)=\xi$. Поскольку
$\pi$ --- локальный гомеоморфизм (теорема~\ref{thm:localhomeo}), найдётся
открытая окрестность $N\ni\xi$ такая, что $\pi|_{N}$ --- гомеоморфизм
на открытое $V\ni x$. По непрерывности $s'$ и $s''$ множества
$(s')^{-1}(N)$ и $(s'')^{-1}(N)$ открыты и содержат $x$; их пересечение
$W$ открыто, и для любого $y\in W$ оба ростка $s'(y),s''(y)$ лежат в
$N$ и проецируются в $y$. Так как $\pi|_N$ инъективно,
$s'(y)=s''(y)$. Значит $W\subset E$, т.\,е. $E$ открыт.

Шаг 2: по условию $U_i\subset E$ для каждого $i$, а $U=\bigcup_iU_i$,
поэтому $U\subset E$, т.\,е. $s'=s''$ на всём $U$.

Заметьте, что <<закрытость $E$» здесь не нужна: достаточно открытости
$E$ и того, что $U_i\subset E$.
\hfill$\square$

\subsection{Аксиома (F2): склейка}
\textbf{Задача.} $U=\bigcup_iU_i$ и даны
$s_i\in\Shv{F}^{\#}(U_i)$, согласованные:
$s_i|_{U_i\cap U_j}=s_j|_{U_i\cap U_j}$.

\textbf{Доказательство.} Определим $s\colon U\to\et(\Shv{F})$ по правилу
$s(x):=s_i(x)$ для любого $i$ с $x\in U_i$.
\begin{conditions}
\item \emph{Согласованность выбора}: если $x\in U_i\cap U_j$, то
 $s_i(x)=s_j(x)$, потому что сечения $s_i$ и $s_j$ совпадают на
 $U_i\cap U_j$ (это их ограничение одно и то же, а сечение
 определяется значениями). Значит $s$ --- функция.
\item \emph{$\pi\circ s=\id$}: выполняется по построению, ведь
 $s_i(x)\in\Shv{F}_x$.
\item \emph{Непрерывность}: пусть $G\subset\et(\Shv{F})$ открыто. Тогда
 \[
   s^{-1}(G)\cap U_i=s_i^{-1}(G)
 \]
 открыто в $U_i$ (непрерывность $s_i$), а $U_i$ открыто в $U$,
 значит $s^{-1}(G)\cap U_i$ открыто и в $U$. Наконец
 \[
   s^{-1}(G)=\bigcup_i\bigl(s^{-1}(G)\cap U_i\bigr)
 \]
 --- объединение открытых в $U$, значит открыто.
\item \emph{Ограничения}: по построению $s|_{U_i}=s_i$.
\end{conditions}
Уникальность $s$ тривиальна: сечение определяется своими значениями.
\hfill$\square$

Итак, $\Shv{F}^{\#}$ --- действительно пучок (и пучок абелевых групп,
если $\Shv{F}$ --- предпучок групп).

\begin{risunok}{Проверка аксиом для $\Shv{F}^{\#}$: слева --- (F1)
 (равны на каждом $U_i$ $\Rightarrow$ равны), справа --- (F2)
 (согласованные $s_1,s_2$ склеиваются в $s$). \label{fig:axioms}}
\resizebox{\linewidth}{!}{\begin{tikzpicture}[>=Stealth,font=\small]
  % левая панель: (F1)
  \begin{scope}
    \draw[fig base] (0,0) ellipse (1.9 and 0.6);
    \begin{scope}
      \clip (0,0) ellipse (1.9 and 0.6);
      \fill[primary!10] (-0.7,0) ellipse (1.35 and 0.55);
      \fill[success!10] (0.7,0) ellipse (1.35 and 0.55);
    \end{scope}
    \draw[fig open] (-0.7,0) ellipse (1.35 and 0.55);
    \draw[fig open2] (0.7,0) ellipse (1.35 and 0.55);
    \node[primary] at (-1.7,0.8) {$U_1$};
    \node[success] at (1.7,0.8) {$U_2$};
    \draw[fig fiber] (-1.0,0.51) -- (-1.0,2.35);
    \draw[fig fiber] ( 1.0,0.51) -- ( 1.0,2.35);
    \fill (-1.0,1.75) circle (2.6pt);
    \fill ( 1.0,1.95) circle (2.6pt);
    \draw[fig sect,smooth] plot coordinates
      {(-2.1,1.4) (-1.0,1.75) (0,1.5) (1.0,1.95) (2.1,1.7)};
    \draw[danger,line width=1pt,dashed,smooth] plot coordinates
      {(-2.1,1.75) (-1.0,1.75) (0,2.05) (1.0,1.95) (2.1,2.15)};
    \node[danger] at (-2.45,1.2) {$s'$};
    \node[danger,font=\small] at (2.45,2.3) {$s''$};
    \node[font=\footnotesize] at (0,-1.05)
      {(F1): равны на каждом $U_i$ $\Rightarrow$ равны};
  \end{scope}
  % правая панель: (F2)
  \begin{scope}[xshift=6.6cm]
    \draw[fig base] (0,0) ellipse (2.4 and 0.7);
    \begin{scope}
      \clip (0,0) ellipse (2.4 and 0.7);
      \fill[primary!10] (-0.85,0) ellipse (1.55 and 0.65);
      \fill[success!10] (0.85,0) ellipse (1.55 and 0.65);
    \end{scope}
    \draw[fig open] (-0.85,0) ellipse (1.55 and 0.65);
    \draw[fig open2] (0.85,0) ellipse (1.55 and 0.65);
    \node[primary] at (-2.2,0.85) {$U_1$};
    \node[success] at (2.2,0.85) {$U_2$};
    \draw[fig fiber] (-1.8,0.463) -- (-1.8,2.65);
    \draw[fig fiber] ( 0,0.7)     -- ( 0,2.65);
    \draw[fig fiber] ( 1.8,0.463) -- ( 1.8,2.65);
    \fill (-1.8,1.45) circle (2.6pt);
    \fill ( 0,2.15)   circle (2.6pt);
    \fill ( 1.8,1.65) circle (2.6pt);
    \draw[fig leafA,smooth] plot coordinates
      {(-2.4,1.3) (-1.8,1.45) (-0.9,1.85) (0,2.15)};
    \draw[fig leafB,smooth] plot coordinates
      {(0,2.15) (0.9,2.0) (1.8,1.65) (2.4,1.5)};
    \node[primary] at (-1.3,1.22) {$s_1$};
    \node[success] at (1.25,1.28) {$s_2$};
    \node[danger] at (0.42,2.52) {$s$};
    \node[font=\footnotesize] at (0,-1.15)
      {(F2): $s_1,s_2$ согласованы $\Rightarrow$ есть $s$};
  \end{scope}
\end{tikzpicture}}
\end{risunok}

Пояснение к рисунку~\ref{fig:axioms}. Слева --- аксиома (F1):
две кривые $s'$ (сплошная) и $s''$ (пунктирная) заранее разные,
но над каждым $U_i$ их значения совпадают --- это и есть общие
точки на волокнах; (F1) утверждает, что тогда $s'=s''$
отождествляются полностью. Справа --- аксиома (F2): локальные
сечения $s_1$ (синее, над $U_1$) и $s_2$ (зелёное, над $U_2$)
согласованы в общей части $U_1\cap U_2$ --- обе кривые проходят
через одну и ту же точку на общем волокне, --- и их склейка
даёт глобальное сечение $s$.

\sectionsrc{photo\_2 (стр.~10): <<внезапная вставка с F1 и F2
 для $\et(\Shv{F})$» --- с рисунками, на которых листы $\Shv{F}^{\#}$
 над покрытием $U=U_i\cup U_j$ показаны <<зелёными и красными»;
 отмечено, что мало проверить равенство образов, надо именно равенство
 сечений.}

\begin{bookbox}
\booksrc{Годеман, п.~3.9, с.~183---184: каноническое отображение
 $\Shv{F}(X)\to\Shv{F}^{\#}(X)$ инъективно, каково бы ни было
 $\Shv{F}$, если выполняется (F1); если для $\Shv{F}$ выполнена (F2)
 и $X$ паракомпактно, то это отображение --- на всё множество.
 Общая идея: в $\Shv{F}(X)$ <<локально равные» сечения отождествляются,
 а в $\Shv{F}^{\#}(X)$ <<добавляются» сечения, склеенные из
 согласованных локальных наборов.
}
\end{bookbox}

%==================================================================
